\begin{tcolorbox}[colback=gray!5,colframe=black,title=\natcap{} Annotation Prompt,fonttitle=\bfseries]
You are an image annotator expert.\\\\
You will receive an image along with its classification label and the classification task scope, and your task is to provide contextualized metadata about it.\\\\
The output should be a JSON object with the following metadata fields:
\begin{itemize}
  \item \textbf{caption}: A descriptive caption of the image accounting for its label. This should be a \textbf{unique} and concise sentence that describes the image in detail.
  \item \textbf{class\_tags}: A list of tags that represents the image and can help identify the class. (e.g., for a car image with its model as a class, this could be some specific attribute of the car)
  \item \textbf{other\_tags}: A list of tags that represents the image but can help identify the image among others of the same class. (e.g., for a car image with its model as a class, this could be its color or the background of the image)
  \item \textbf{is\_image\_class\_explicit}: Boolean, could the class be inferred from the image alone? (e.g., the class is a country and you cannot necessarily infer it from the image alone, so this would be \texttt{false})
\end{itemize}

Please ensure that the output is in valid JSON format.\\

\textbf{Example:}\\
You receive an image of what is clearly a car with its model as a class (here Audi TTS coupe 2012) for a car model classification task.\\
The output could be a JSON object like this:

\begin{lstlisting}[basicstyle=\ttfamily\small]
{
  "caption": "A red Audi TTS coupe 2012 car parked on a sunny street 
                in front of a sport shop.",
  "class_tags": ["sport coupe","four door coupe","17'' alloy wheels"],
  "other_tags": ["sunny street","parked","red","sport shop"],
  "is_image_class_explicit": true
}
\end{lstlisting}

\textbf{------}\\
Classification scope: \{task\_info\}\\
Image label: \{label\}\\
Answer:
\label{box:natcap_prompt}
\end{tcolorbox}